\documentclass{article}
\usepackage{graphicx}
\usepackage{amsmath}
\usepackage{amsfonts}
\usepackage{bookmark}
\usepackage{booktabs}
\usepackage{exercises}
\usepackage{csquotes}
\graphicspath{ {./images/4.2} }
\usepackage{indentfirst}
\usepackage{amsthm}
\usepackage{yhmath}
\usepackage{placeins}
\newcommand*{\oldneg}{\mathord{\sim}}
\newcommand\showdiv[1]{\overline{\smash{)}#1}}
\newcommand\Mydiv[1]{%
\smash{\raise.3ex\hbox{$\big)$}}\mkern-8mu
        \overline{\enspace\strut#1}}
\setlength{\tabcolsep}{5pt}
\theoremstyle{definition}
\newtheorem{definition}{Definition}

\setlength{\parindent}{.5in}
\setlength{\parskip}{12pt}
\renewcommand{\arraystretch}{2}
\newcommand{\?}{\stackrel{?}{=}}

\title{Extension Answers (pg. 242)}
\begin{document}
\maketitle

\section*{1}
We are given the model \[ y=750(1.02)^t \] where $t$ is measured in weeks. We wish to find the daily percent increase. To do this, we can use the fact that 
7 days are in a week, and $\frac{1}{7}(7t)=t$: \[ y=750(1.02)^{\frac{1}{7}(7t)} \] \[ =750(1.02^\frac{1}{7})^{7t}=750(1.00283)^7t \]
meaning the daily increase is about $0.00283$, or about $0.283\%$.


I was a bit confused why the book didn't present this as a composition of functions. Instead of saying \enquote{we can use the fact that 7 days are in a week, 
and $\frac{1}{7}(7t)=t$}, we could simply let $d=\frac{1}{7}t$ (because a day is one seventh of a week), and then find $f(d(t))$ which gives:
\[ y=750(1.02)^\frac{1}{7}t \] and then we get the same result (by separating $\frac{1}{7}$ and $t$), but this time really we really are measruing just one day. 

\section*{2}
We are given the model \[ A=n(\frac{1}{2})^\frac{t}{15} \]  where $t$ is measured in hours. We wish to find the hourly decay rate, which can be done by 
using the fact that $t=\frac{1}{15}t$: \[ A=n(\frac{1}{2})^{\frac{1}{15}t} \] \[ A=n((\frac{1}{2})^\frac{1}{15})^t=n(0.95484)^t \] meaning the hourly decrease is 
about $1-0.95484=0.04516$, or about $4.516\%$.

\section*{6}
We are given that the population doubles every 3 days, and are asked for the daily percent increase. Let $P$ be the principal population, and $A$ be the population after 
$t$ days. Then we know that:
\[ A=P\times 2^1=P\times 2^\frac{3}{3},when\; t=3 \] \[ A=P\times 2^2=P\times 2^\frac{6}{3},when\; t=6 \] \[ A=P\times 2^3=P\times 2^\frac{9}{3},when\; t=9 \]
meaning a model to represent this scenario is \[ A=P(2)^\frac{t}{3} \] Next, we find the daily percent increase:
\[ A=P(2)^\frac{t}{3} \] \[ =P(2^\frac{1}{3})^t \] \[ =P(1.26)^t \] meaning the daily percent increase is about $0.26$, or about $26\%$.

\end{document}
